% =====================================================================
% Reassessment of the progress made on reducing malaria mortality in Africa
% Neutral preprint format (medRxiv-style); structured for a general
% medical journal (Lancet/NEJM): structured abstract, IMRaD, Vancouver refs.
% Compiles with pdflatex. Target main-text length ~3200 words.
% =====================================================================
\documentclass[12pt]{article}

\usepackage[utf8]{inputenc}
\usepackage[T1]{fontenc}
\usepackage{lmodern}
\usepackage[margin=1in]{geometry}
\usepackage{setspace}
\usepackage{graphicx}
\usepackage{booktabs}
\usepackage{amsmath}
\usepackage{caption}
\usepackage[hidelinks]{hyperref}
\usepackage{lineno}

\captionsetup{labelfont=bf,font=small}
\setstretch{1.5}
\linenumbers

\title{Reassessing the malaria mortality burden in African children}

\author{
James A. Watson$^{1,2,3,*}$, Dhruv Darji$^{1,2}$, [...], Philippe J Guerin$^{1,2}$, Robert W Snow$^{2,4}$
}

\date{\today}

\begin{document}
\maketitle


\small\noindent\textbf{1} Infectious Diseases Data Observatory, Oxford, UK;\\
\small\textbf{2} Center for Global Health and Tropical Medicine, University of Oxford;\\
\small\textbf{3} GiveWell, San Francisco, USA;\\
\small\textbf{4} Population and Health Impact Surveillance Group, KEMRI-Wellcome Trust Research Programme, Nairobi, Kenya;\\
\small $^{*}$Correspondence: jwatowatson@gmail.com





% 300 word limit

\begin{abstract}
    Malaria mortality in African children is estimated from complex geospatial models taking as primary input verbal autopsy data. Verbal autopsy lacks specificity and sensitivity in identifying malaria death, cannot capture indirect mortality effects, and assumes each death has a single cause. We propose a simpler framework for estimating the total burden of malaria attributable mortality in children under 5 years by framing transmission intensity (measured by prevalence of infection in children: \textit{Pf}PR$_{2-10}$) as a risk factor for all-cause death.
    We pooled all available survey data from sub-Saharan Africa since 2000 that recorded complete birth histories and potential confounding variables (136 surveys, 37 countries). Individual child trajectories, discretised into 7 age bands from birth until 5 years, were merged with year and location matched population weighted \textit{Pf}PR$_{2-10}$ estimates from the Malaria Atlas Project. After adjusting for plausible confounding variables, a robust age-dependent association between \textit{Pf}PR$_{2-10}$ and all-cause mortality was observed in data from 2,335,868 children under 5 of whom 105,114 died. \textit{Pf}PR$_{2-10}$ only weakly predicted neonatal death, but was a strong predictor of death from six months and older. At 20\% prevalence, the model attributes 4\% of neonatal deaths, 25\% of deaths at 6-11 months and 43-55\% of deaths at 1-4 years to malaria.
    Applied to national all-cause mortality estimates from 42 sub-Saharan African countries, we estimate that under-5 malaria mortality fell from 10.1 to 4.3 deaths per 1,000 child-years between 2000 and 2015 (a 57\% decline), and then to 3.3 deaths per 1,000 child-years in 2024 (a further 24\% decline). IHME and UN IGME estimate similarly large declines up until 2015 (49 and 60\% decrease, respectively) but little change since (8 and 6\% decrease). We estimate malaria mortality rates in 2024 were 37\% and 33\% higher than IHME and UN IGME estimate, respectively. Recent progress against childhood malaria has been substantially greater than currently reported, yet the malaria-attributable burden of child mortality remains larger than current cause-of-death estimates suggest.
\end{abstract}

% =====================================================================
\section*{Introduction}


%Segue into a) malaria infection is likely to have significant indirect effects on child survival and b) infection prevalence and U5M are measured more often, at wider geographic scales and are far less ambiguous outcomes.

%Why is this important now in climate of reduced funding etc....
%Here we...... I think a reader would benefit from a short punchy statement of what analysis you present in introduction. i.e. if u use MA fractional mortality vs all-cause mortality is temporal impact claimed bigger/smaller or the same.

% How do we estimate malaria mortality?
Most contemporary publications report that around half a million preventable deaths in children in sub-Saharan Africa each year  are attributable to malaria  \cite{wmr}. The highest burden of disease and death occurs in difficult to reach and poorly resourced areas where access to prevention and treatment is limited. Knowing the number of deaths caused by malaria and where they are located is critical for guiding allocation of scarce resources, especially now with reduced overseas development aid for health. Over the last 25 years, substantial investments have been made in modelling the burden of malaria mortality \cite{gates2017ihme}. The most influential and often cited estimates are from the World Health Organization (WHO) yearly reports, the openly available model based predictions from the Institute for Health Metrics and Evaluation (IHME) and the United Nations Inter-agency Group for Child Mortality Estimation (UN IGME). Their methodologies differ, but the primary input data are mostly the same. These models all assume that every death has a single cause, the proportions of which are primarily estimated from verbal autopsy (VA) studies. For malaria, estimating malaria mortality from VA data is problematic because it ignores well characterised indirect mortality effects \cite{scott,Uyoga:19,Muriuki:21} and because VA studies cannot reliably attribute causality to malaria versus other infectious diseases \cite{todd1994,snowva1992}. Deaths caused by malaria are notoriously difficult to ascertain in the absence of supporting parasitological confirmation, either from microscopy or a rapid diagnostic test \cite{reyburn2004overdiagnosis}. The symptoms of severe disease are non specific and cannot be easily differentiated from other febrile illnesses \cite{white_misdiagnosis,Watson:22}. In addition, VA studies from representative samples are infrequently conducted, data derive from only a few demographic surveillance sites in sub-Saharan Africa and mostly during times of failing first line treatments [REF].

% Where is the burden
In most of malaria endemic sub-Saharan Africa, malaria is primarily a childhood disease. Entomological inoculation rates - the number of infectious bites received each year - can be in the tens to hundreds. Acquisition of disease controlling immunity follows multiple parasite exposures from birth \cite{filipe2007immunity} whereby most infections do not result in symptomatic disease. Parasite controlling immunity develops even more slowly, so most older children are usually asymptomatically infected. Transmission intensity, a proxy for inoculation rates, can be tracked by the proportion of children who are infected at any given point, most meaningfully when restricted to children between 2 and 10 years of age (\textit{Pf}PR$_{2-10}$) \cite{Metselaar:59,smith2007standardizing,smith2005entomological}. In areas of high transmission (e.g. \textit{Pf}PR$_{2-10}$>20\%), very few infections develop into life-threatening severe malaria \cite{duffy}, and these are concentrated in the first few years of life \cite{Kamau:22}. Severe malaria represents an inability of the child to control parasite multiplication and has high acute mortality \cite{white_severe}. However, repeated infections and chronic parasite carriage have other indirect negative consequences. Most hospital admissions due to malaria in high transmission settings are for severe anaemia \cite{Kamau:22}. In addition to chronic anaemia, malaria increases the risk of bacterial sepsis \cite{scott,Mackenzie:10}, iron deficiency \cite{Muriuki:21}, and reduces growth \cite{Oldenburg:18}. There is a clear relationship between transmission intensity measured by \textit{Pf}PR$_{2-10}$ and admission to hospital for severe malarial anaemia, and severe malaria more generally \cite{snowsevere,Paton2021,Kamau:22}. If \textit{Pf}PR$_{2-10}$ can predict the incidence of severe forms of malaria, and if indirect effects are a fixed proportion of the direct effects, then it follows that \textit{Pf}PR$_{2-10}$ should predict the overall number of direct and indirect deaths due to malaria. Because prevalence of infection and all-cause childhood mortality rates are measured at regular intervals at a wide geographic scale from representative surveys, these are the most reliable inputs for estimating the total burden of disease. 

We propose an alternative approach to estimating the number of deaths caused by malaria mortality, combining direct and indirect effects. Rather than trying to estimate the number of malaria deaths directly, we propose a framework in which the intensity of malaria transmission is a risk factor for death, analogous to how the mortality burden of malnutrition is determined \cite{ritchie:24}. This alternative approach estimates the relationship between \textit{Pf}PR$_{2-10}$ and age-specific all-cause mortality, adjusting for the known determinants of childhood death that could confound this relationship. We show that the relationship between malaria transmission intensity as measured by \textit{Pf}PR$_{2-10}$ and age-specific all-cause mortality is extremely robust based on 25 years of representative household surveys in Africa. These data suggest that malaria causes more deaths in children under 5 than currently estimated, but that the progress in reducing mortality in the last 2 decades has been greater than previously reported.


\section*{Results}
\subsection*{\textit{Pf}PR$_{2-10}$ and all-cause mortality}

We pooled all available microdata from Demographic and Health Surveys (DHS), Malaria Indicator Surveys (MIS), and Multiple Indicator Cluster Surveys (MICS) conducted in main land sub-Saharan Africa or Madagascar since the year 2000 which recorded complete birth histories. This gave a total of 166 eligible surveys (117 DHS, 5 MIS, 44 MICS) across 39 countries (Figure \ref{fig:survey_timing}). Each child trajectory from birth was divided into discrete age bands of $<$1 month (neonatal period), 1-5 months, 6-11 months, and then yearly, corresponding to how data are usually recorded in the complete birth histories. An age band was included in the analysis if the child entered it within the five years before the interview and the band ended (i.e. was completed) before the interview. Each child trajectory was then matched to a sub-national malaria-endemic area (taken from \cite{snow:17}, mostly first-level administrative units) and matched with population-weighted \textit{Pf}PR$_{2-10}$ estimate from the Malaria Atlas Project (MAP) corresponding to the calendar year the band started \cite{map2025}, along with a series of national or subnational estimates of identified confounding variables. The five surveys from Lesotho (four DHS, one MICS) were excluded because Lesotho is not a malaria endemic country. 

Restricting the analysis to areas and years with all 17 confounding variables measured, a total of 136 surveys (99 DHS, 37 MICS) in 37 countries were included in the primary analysis, representing the complete birth histories of 1,125,610 mothers (Figures \ref{fig:map} and \ref{fig:flow}); all five MIS surveys, 14 DHS surveys and 6 MICS surveys were excluded because of missing data on identified confounders. The included surveys contributed 7,639,443 child age band records from a total of 2,335,868 children with 105,114 deaths recorded, in 1,452 survey-by-area units. 37\% of these deaths were in neonates (first month of life), 66\% in the first year of life, and 13\% between 1 and 2 years of age.


\begin{figure}
    \centering
    \includegraphics[width=0.95\linewidth]{figures/fig1_analysis_units_map.png}
    \caption{Number of complete birth histories included in the primary analysis by administrative unit used as the unit of analysis. Administrative units are based on \cite{snow:17}.}
    \label{fig:map}
\end{figure}

We fitted separate models for each age band to estimate the relationship between \textit{Pf}PR$_{2-10}$ and all-cause age-specific mortality, adjusted for a set of 17 administrative area-level and national covariates that could plausibly confound this relationship (socio-economic, health-system, nutritional, governance and child HIV incidence measures). The estimated relationship between \textit{Pf}PR$_{2-10}$ and all-cause mortality varied substantially across age groups (Figure \ref{fig:main_result} and Table \ref{tab:primary-age-band-effects}). For neonatal mortality (37\% of included deaths), there was only a weak relationship with \textit{Pf}PR$_{2-10}$. Neonates are protected against malaria due to passive immunity from maternal antibodies and a high proportion of foetal haemoglobin \cite{WHO:14}. Most neonatal deaths in these areas are caused by birth-related complications (preterm birth, birth asphyxia or trauma), or bacterial and viral infections and congenital anomalies \cite{who_newborn_mortality_2024}. Malaria in pregnancy increases the risk of low birth weight which is a predictor of neonatal and infant death but is thought to only contribute to a small proportion of neonatal death \cite{Walker:14}. As such, even in high malaria transmission areas, we would not expect a large proportion of neonatal deaths to be caused by malaria. The weak observed association between \textit{Pf}PR$_{2-10}$ and neonatal death suggests that there is no major residual confounding. In all other age groups however, there was a clear relationship between \textit{Pf}PR$_{2-10}$ and all-cause mortality. Post-neonatal infant mortality (death between 1 and 11 months) increased monotonically with increasing \textit{Pf}PR$_{2-10}$. 
In older children, however, the increase in risk slowed markedly above prevalences of around 30\%, and for children aged 3 and 4 years the curves were essentially flat above prevalences of $\sim$40\%, but with large uncertainty in the estimates given there were less data in these ranges. This pattern fits our understanding of the epidemiology of malaria in very high transmission settings whereby disease is concentrated in first years of life \cite{Kamau:22}, and therefore older children who have not died from malaria by age 3 or 4 have disease controlling immunity.

Under this model, for an average \textit{Pf}PR$_{2-10}$ of 10\%, approximately 30\% of all deaths in children 12 months and older were estimated to be caused by malaria (28–37\% across the four yearly bands), compared with less than 5\% in neonates (Figure \ref{fig:main_result2}). For an average \textit{Pf}PR$_{2-10}$ of 40\% (upper end of prevalence estimates in sub-Saharan Africa now), there was a sharper age-dependent increase in malaria deaths, with around half of all deaths between 12 and 59 months attributable to malaria (47–59\%), peaking at 59\% between 24–35 months.

\begin{figure}[htbp]\centering
\includegraphics[width=0.95\linewidth]{figures/fig2_malaria_attributable_fraction_by_age.png}
\caption{Estimated malaria-attributable share of all-cause mortality within each age-band at different prevalences of infection (\textit{Pf}PR$_{2-10}$). Points are point estimates and vertical bars 95\% confidence intervals; lines connect the discrete age bands.}
\label{fig:main_result2}
\end{figure}

We conducted a series of sensitivity analyses in subsets of the primary analysis dataset (see Methods): survey regions in the Sahel region only (31 surveys in 8 countries), where transmission is highly seasonal; each UN sub-region of Africa (Western Africa, 61 surveys in 15 countries; Central Africa, 18 surveys in 7 countries; Eastern Africa, 50 surveys in 12 countries; Southern Africa was not refitted because its \textit{Pf}PR$_{2-10}$ never exceeds 10.3\%); surveys conducted up to (74 surveys) or after (62 surveys) the median survey year, 2014; and child age bands entered up to or after the median calendar date of entry (late 2012). In general the subset fits gave similar results (Figures \ref{fig:subgroups}–\ref{fig:subgroups_period}). The exceptions were the Sahel region, where mortality rose more steeply between prevalences of 20\% and 40\% for children between 6 and 59 months, and Central Africa, where mortality did not increase between prevalences of 20\% and 40\% after the first month. In the unadjusted data (Figure \ref{fig:survival_by_region}), most of Central Africa's records at low prevalence come from Chad (58\% of its records at 24–59 months with a prevalence below 10\%), where survival at these ages (99.1\%) is lower than at similar prevalence elsewhere (99.6\%). A sensitivity analysis using multiple imputation for the missing covariates (8,630,306 child age-band records and 122,092 deaths) gave very similar results to the primary analysis restricted to complete cases (Figure \ref{fig:MI}). Because acute or chronic malnutrition may mediate the relationship between malaria transmission and death as well as confound it, we refitted the models on the same sample without adjusting for wasting and stunting; the estimated prevalence-mortality relationship was very similar (hazard ratios changed by at most 0.04; Figure \ref{fig:without_malnutrition}).

\subsection*{Comparison with IHME and UN IGME malaria mortality estimates}

If we interpret the estimated relationships between \textit{Pf}PR$_{2-10}$ and age-specific all-cause mortality as causal (i.e. there are no unmeasured confounders), and thus as generalisable across geographies in sub-Saharan Africa, then the number of deaths caused by malaria (directly and indirectly) can be estimated by applying the age band specific hazard ratios contrasting the observed prevalence with zero (or any arbitrary very low prevalence at which malaria deaths are assumed to be negligible) to national all-cause death rates. We applied the fitted age band effects to IHME all-cause deaths by age band and population-weighted national \textit{Pf}PR$_{2-10}$ in all 42 sub-Saharan African countries with MAP predicted prevalence estimates, for the year 2024, and compared the results with the cause-specific under-5 malaria death estimates of IHME and of UN IGME, expressed per 1,000 under-5 child-years (Figure~\ref{fig:IHME_v_model}A).

Pooled across the 42 countries, our estimated under-5 malaria mortality rate fell from 10.1 to 4.3 deaths per 1,000 child-years between 2000 and 2015 (a 57\% decline) and to 3.3 in 2024 (a further 24\% decline), whereas the IHME and UN IGME rates fell by 49\% and 60\% from 2000 to 2015 but by only 8\% and 6\% thereafter (Table~\ref{tab:source-comparison}; Figure~\ref{fig:IHME_v_model}C). In 2024 the \textit{Pf}PR ACM model eestimates approximately 587,000 under-5 malaria deaths in these countries, 37\% more than IHME (428,000) and 33\% more than UN IGME (440,000).


Although the country specific under 5 mortality rates were highly correlated between IHME and our proposed new approach (\textit{Pf}PR ACM, Spearman $\rho$=0.91), the \textit{Pf}PR ACM model estimate exceeded IHME in 34 of the 42 countries (median ratio 1.5), with the largest absolute differences in several populous, high-burden countries (Table \ref{tab:country}).
In Nigeria, IHME estimates approximately 138 thousand deaths and UN IGME approximately 165 thousand deaths, whereas our approach estimates 207 thousand deaths (50\% and 25\% higher, respectively).
In Chad we estimate a 3.8-fold higher total burden than IHME and 1.3-fold higher than UN IGME. Within Nigeria, applying the \textit{Pf}PR ACM model to state-level age-specific all-cause mortality estimates (IHME) and state-level \textit{Pf}PR$_{2-10}$ estimates gives malaria mortality estimates greater than IHME's estimates in 35 of the 37 states (Figure \ref{fig:IHME_v_model}B). The differences are concentrated in the northern states: the median state ratio is 1.7 in the three northern zones and 1.2 in the three southern zones. Expressed as the probability of dying before age 5 (Figure \ref{fig:cumulative_risk}), the model attributes a median of 12 deaths per 1,000 live births to malaria across the 42 countries in 2024 (a median 20\% of the all-cause probability), with the highest values in South Sudan (33 per 1,000), Sierra Leone (31) and the Central African Republic (31); this is a median 1.5 times the probability implied by IHME's cause-specific malaria death rates.


\begin{figure}[htbp]\centering
\includegraphics[width=\linewidth]{figures/fig3_burden_comparison.png}
\caption{Comparison of national and subnational (Nigeria only) estimates of malaria mortality rates in children under 5 between the PfPR-ACM model and IHME and UN IGME. A: Comparison at national level; B: Comparison at the subnational (state) level for Nigeria; C: Annual predicted malaria mortality rates per 1,000 child years for children under 5 in sub-Saharan Africa since 2000.}
\label{fig:IHME_v_model}
\end{figure}

\begin{table}[htbp]
\centering
\caption{Malaria mortality by source in 2000, 2015 and 2024, and change to 2024. Negative values are reductions.}
\label{tab:source-comparison}
\small
\setlength{\tabcolsep}{3pt}
\renewcommand{\arraystretch}{1.15}
\begin{tabular*}{\linewidth}{@{\extracolsep{\fill}}lrrr@{}}
\toprule
& \shortstack{PfPR-ACM\\model} & IHME & \shortstack{UN\\IGME} \\
\midrule
\multicolumn{4}{@{}l}{\textit{Malaria deaths}} \\
\quad 2000 & 1,111,970 & 563,657 & 724,276 \\
\quad 2015 & 687,031 & 416,122 & 416,815 \\
\quad 2024 & 586,723 & 428,147 & 440,123 \\
\quad Change 2000--2024 & $-47.2\%$ & $-24.0\%$ & $-39.2\%$ \\
\quad Change 2015--2024 & $-14.6\%$ & $+2.9\%$ & $+5.6\%$ \\
\addlinespace
\multicolumn{4}{@{}l}{\textit{Malaria mortality rate}} \\
\quad 2000 & 10.12 & 5.13 & 6.59 \\
\quad 2015 & 4.33 & 2.62 & 2.63 \\
\quad 2024 & 3.31 & 2.41 & 2.48 \\
\quad Change 2000--2024 & $-67.3\%$ & $-53.0\%$ & $-62.4\%$ \\
\quad Change 2015--2024 & $-23.7\%$ & $-8.1\%$ & $-5.7\%$ \\
\bottomrule
\end{tabular*}
\par\vspace{0.5em}
\begin{minipage}{\linewidth}
\footnotesize
All three columns are deaths before age five in the same 42 countries. Rates use the same annual under-five person-years, implied by the IHME all-cause death count and rate, so within a year the three sources share one denominator.
Model deaths are IHME all-cause deaths in each age band multiplied by the estimated malaria-attributable fraction; IHME and UN IGME are cause-specific malaria death estimates. These are different estimands.
For IHME and UN IGME, deaths rise but rates fall after 2015, because under-five person-time in these countries grew by 12\% between 2015 and 2024.
\end{minipage}
\end{table}


\begin{table}[htbp]
\centering
\caption{Ten countries with the largest absolute differences between the malaria-attributable deaths of the PfPR-ACM model and the malaria deaths estimated by IHME, before age five in 2024.}
\label{tab:country}
\begin{tabular}{lrrr}
\toprule
Country & PfPR-ACM model & IHME & UN IGME \\
\midrule
Nigeria & 206,863 & 138,111 & 165,329 \\
Niger & 36,238 & 20,343 & 23,877 \\
Chad & 18,646 & 4,884 & 14,345 \\
Mozambique & 17,971 & 10,285 & 10,315 \\
Mali & 19,629 & 12,945 & 6,984 \\
South Sudan & 10,635 & 4,779 & 3,980 \\
Cameroon & 21,413 & 16,295 & 6,286 \\
Angola & 19,684 & 14,652 & 4,734 \\
Central African Republic & 8,629 & 3,865 & 5,220 \\
Cote d'Ivoire & 17,098 & 12,848 & 5,753 \\
\midrule
All 42 countries & 586,723 & 428,147 & 440,123 \\
\bottomrule
\end{tabular}
\par\smallskip
\begin{minipage}{\linewidth}\footnotesize
UN IGME denotes the direct under-five CA-CODE 2026 series. Totals include all 42 estimable countries, not only the ten shown. Death counts are rounded after estimation. The model uses the MAP fits with $\gamma=2$ on Snow et al. polygon units.
\end{minipage}
\end{table}



\begin{figure}
    \centering
    \includegraphics[width=0.95\linewidth]{figures/fig4_under5_death_probability.png}
    \caption{National \textit{Pf}PR$_{2-10}$ (left) and probability of death by 5 years of age attributable to malaria versus all causes across 42 sub-Saharan African countries (right).}
    \label{fig:cumulative_risk}
\end{figure}


\section*{Discussion}

%% Progress in the last quarter century
The current dominant narrative regarding the global burden of malaria mortality is that progress has stalled since 2015 and that the total number of malaria deaths (predominantly in young children) has plateaued at approximately 600,000 per year \cite{wmr,gbd}. 
Compared with the year 2000 - a pivotal year in malaria control history with the start of the Roll Back Malaria Partnership \cite{doi:10.1126/science.280.5372.2067} - malaria mortality estimates for 2024 from IHME and UN IGME in sub-Saharan Africa represent decreases of 53\% and 62\%, respectively, in the number of deaths per 1,000 child-years in children under 5. A quarter of a century ago chloroquine was the main available antimalarial with high-grade resistance ubiquitous across the continent, and very few households had access to long lasting insecticide treated bednets \cite{white_averting}. A decrease of less than two-thirds over a quarter century represents a failure of the Roll Back Malaria campaign, which aimed to halve mortality by 2010 and then halve again by 2015 \cite{doi:10.1126/science.280.5372.2067}. This narrative of disappointing progress has recently been challenged: why is it that the situation has changed so little after so much funding for malaria control \cite{white:22}? In comparison, our PfPR-ACM model estimates a larger 67\% decrease between 2000 and 2024 (57\% by 2015, against 49\% for IHME and 60\% for UN IGME; the RBM goal was a 75\% reduction by 2015) and continued declines since 2015: a 24\% decrease between 2015 and 2024, against 8\% (IHME) and 6\% (UN IGME) (Table \ref{tab:source-comparison}).

IHME and UN IGME differ in how they estimate malaria mortality, but both depend on layering multiple complex models. In both cases, the key inputs are prevalence (\textit{Pf}PR$_{2-10}$) estimates and verbal autopsy data determining the share of all-cause deaths attributable to malaria. The IHME algorithm converts prevalence into incidence of clinical disease using an ensemble of transmission models. The proportion of untreated clinical cases is then estimated from survey data asking about treatment seeking behaviour, appropriate drug use and drug efficacy. The untreated case fatality ratio is then estimated by multiplying the share of malaria deaths (based on verbal autopsy) by the all-cause mortality divided by the untreated incidence \cite{ihme}. The approach taken by UN IGME is simpler. The share of all-cause deaths in children under 5 attributable to malaria is estimated from a Bayesian multinomial model fitted to verbal autopsy studies with prevalence as a covariate, in addition to other key variables such as vaccine coverage, overall under 5 mortality rate, calendar year. These models can disagree substantially on the proportion of deaths attributable to malaria at the national level (Supplementary Figure \ref{fig:malaria_share_pfpr}). In both cases, the critical input is verbal autopsy estimates. Verbal autopsy cannot distinguish death due to malaria from death due to other important infections such as bacterial sepsis \cite{valecha:11}. 


One possible reason for the major difference in the proportion of deaths attributable to malaria between this simpler approach and IHME and UN IGME is that verbal autopsy cannot accurately characterise the indirect effect of repeated malaria infection. Evidence from multiple sources, including Mendelian randomisation studies using sickle cell trait (HbAS), points towards malaria having a substantial indirect effect via increased susceptibility to other infections \cite{Uyoga:19,Muriuki:21}. In coastal Kenya, when malaria transmission declined dramatically, the incidence of life-threatening gram-negative bacteraemia was halved \cite{scott}. This indirect effect could amount to around a quarter of all deaths caused by malaria if we take the difference between our estimates and IHME/UN IGME at face value. This suggests that malaria should be considered as a ``risk factor'' similar to malnutrition \cite{Madewell2024}. The exact causal interpretation of the effect of malnutrition on childhood death is unclear because of the ambiguity in how to modify it: low birth weight is associated with around 1.5 million under-5 deaths \cite{ritchie:24}, but low birth weight is unlikely to be the direct cause of death, rather a proxy measurement for other direct causes. In the case of malaria, however, there is a clearer causal interpretation because community prevalence is directly modifiable via interventions such as insecticide treated bednets or indoor residual spraying. Historical data from large scale efforts to reduce or interrupt transmission supports this prediction. For example, at Pare-Taveta, on the Kenyan-Tanzanian border, 3 years of house spraying reduced \textit{Pf}PR$_{2-10}$ from 60\% to 5\% and resulted in a halving of infant and under-five mortality that reverted to pre-intervention levels after spraying ceased in 1958 \cite{Smith1959}.
%Cause shares for ages 1–59 months in high-mortality countries come from a Bayesian multinomial model fitted to verbal-autopsy studies. There are 411 data points from 42 countries, and since 2026 they are calibrated against CHAMPS tissue-sampling data

% Appendix in IHME:
%"Pixel-year predictions of CFR were then multiplied by the untreated incidence rate rasters from the MAP cube to yield pixel-year mortality rate estimates, which were then multiplied by pixel-year population to derive pixel-year malaria death counts."


%Model (Africa):
%MAP maps PfPR2–10 from household surveys and converts it to case incidence using an ensemble of three malaria transmission models.
%Cases not effectively treated are estimated from treatment-seeking, drug-type and drug-efficacy data. This gives "untreated incidence".
%A case-fatality rate among untreated cases is estimated from mostly verbal-autopsy data: the verbal-autopsy malaria share of deaths × national all-cause mortality, divided by untreated incidence. It is then smoothed geographically, adjusting for travel time to cities, age structure, all-cause mortality and sickle-cell prevalence.
%Deaths = case-fatality rate × untreated incidence × population.
%The split into under-5 and older ages comes from separate models fitted to verbal-autopsy and vital-registration age patterns.
%All causes are rescaled so they add up to GBD's own all-cause mortality totals.


%Risk factor - analoguous to malnutrition

%Although it is clear that malaria is a major cause of childhood death in sub-Saharan Africa, it is very difficult to predict the health consequences of elimination or near elimination in all endemic countries. Eradication of malaria within one generation is considered feasible \cite{feachem2019malaria}. Our simple approach for characterising the relationship between transmission intensity and all-cause childhood death provides an answer to how much childhood death could be reduced if malaria was eliminated. In the DRC alone, where prevalence of infection is above 30\%, we estimate that one third of post neonatal childhood deaths could be averted from malaria elimination. 

%Try and highlight more why this is a simple model of impact and one that depends less on VA, imperfect interpolations of sparse data and other variables in your model should be tracked anyway. Prevalence and CEB mortality estimates can be collected at scale with the right funding. etc


%Limitations....
The major limitation of this analysis is that it relies on a causal interpretation of the association between \textit{PfPR}$_{2-10}$ and all-cause mortality. More than half of the variation in prevalence across child records (57\%) is between countries; the rest is within countries, between sub-national areas and over time. However, there is good reason to think that the causal effect is identifiable. The DHS and MICS surveys collect granular data on many important indicators that determine childhood survival, such as access to clean water, vaccination, access to health care, and general socio-economic indicators. After adjustment for these confounding variables, there is only a weak association between \textit{PfPR}$_{2-10}$ and neonatal deaths (37\% of all deaths) but a strong association with older infant and childhood deaths (Figure \ref{fig:main_result}). This striking difference by age group suggests that residual confounding is likely to be minimal. In addition, there is strong biological plausibility for a relationship between prevalence and all-cause death. In many of these sub-Saharan African countries, malaria is thought to be one of the major causes of childhood death. \textit{PfPR}$_{2-10}$ predicts incidence of severe disease \cite{Paton2021}, representing the direct effect of malaria infection on death. It would then follow that \textit{PfPR}$_{2-10}$ should predict all-cause death, reflecting both direct and indirect effects. Prevalence is that of the Snow et al.\ polygon containing most of the population of each record's reported location or, where that location spans several polygons (14\% of mothers), the population-weighted mean over them, which blurs local differences in transmission.
The second main limitation of the work is that it is restricted to children under 5.  Although it is clear that in very high transmission areas, morbidity and mortality from malaria is concentrated in the first years of life \cite{Kamau:22}, individuals still have symptomatic episodes in later childhood and some of these could lead to death. In lower transmission areas, the mortality burden is less concentrated in early childhood \cite{Carneiro:10}. We limited our analysis to children under 5 because this is the most important subpopulation affected by malaria, but also because birth history data are the most reliable for recent deaths in younger children. Age of death is reported at a more granular level for neonates and infants, and 90\% of all reported deaths in the birth histories occurred before age five. How these results might generalise to older children should be considered in future work.

%Discussion needs a powerful end game
Malaria infection impacts child survival substantially more than estimated via its direct effects as a single cause of death. Malaria should be considered as a major risk factor for increased death similar to how the burden of malnutrition is measured \cite{ritchie:24}. Prevalence of infection in children and all-cause under 5 mortality can both be reliably measured via representative household surveys. Our methodology provides a simple framework for countries to evaluate their progress in reducing the burden of malaria mortality without depending on complex models, often considered by national stakeholders as black-boxes. These results predict that reducing malaria parasite exposure will have a larger impact on under-5 mortality than expected from direct malaria deaths alone, similar to what was observed in large scale spraying campaigns done in Africa over 60 years ago.

\section*{Methods}

% Still open after the Snow-unit update (15 October 2026):
% SText 1 describes the archived negative-binomial model and must be rewritten for the binomial cloglog GAM;
% Liberia's child HIV incidence from AIDSinfo counts is described only in the commented-out paragraph below;
% the two different 2024 MAP rasters (national burden vs Snow areas) and the HIV incidence model details are not described;
% Uganda Buganda and Sierra Leone 2013 fixes; licensing of third-party inputs (see docs/PAPER_CODE_FLAGS.md in the code repository).

\subsection*{Data sources}

Complete birth histories were obtained from the Births Recodes of the Demographic and Health Surveys and Malaria Indicator Surveys (DHS/MIS) \cite{dhs} and from UNICEF Multiple Indicator Cluster Surveys (MICS) that recorded a complete birth history \cite{mics}. Birth dates, survival status at the time of the survey interview, and reported age at death were used to construct individual child trajectories. Ages were grouped in 7 age-bands: $<1$ month, 1-5 months, 6-11 months, 1-2 years, 2-3 years, 3-4 years and 4-5 years. 
For DHS/MIS surveys, child, maternal and household characteristics were computed from the survey recodes using the DHS indicator definitions: from the Births Recodes, maternal age at first live birth, maternal education in completed years, household wealth quintile, urban or rural residence, coverage of three doses of DTP-containing vaccine (DTP3) and of measles vaccination among living children aged 12–23 months, the proportion of births in the five years before interview delivered in a health facility, and the proportion of non-first births with a preceding birth interval shorter than 24 months; and, from the Household Member Recodes, household access to an improved drinking-water source, an improved sanitation facility and electricity, and the prevalence of wasting and stunting among children under five years (WHO Child Growth Standards; from the Height and Weight Recodes for surveys whose Household Member Recodes predate the standards). Aggregated to the national level, the eight indicators with published DHS national values reproduced them to within 2 percentage points in 930 of 933 survey–indicator pairs, and in all 790 pairs of the included surveys. Where DTP3 or measles coverage could not be computed for an area, the WHO/UNICEF national estimate (WUENIC) for the survey year was used.
Of the 93 MICS surveys from sub-Saharan Africa, 47 (28 countries) recorded a complete birth history; the other 46 recorded only summary birth histories and thus cannot be used in the age-band analysis. Each MICS birth history was converted to the DHS Births Recode layout and processed with the same age-band construction and eligibility rules as the DHS/MIS data. The 366 MICS region labels were mapped to analysis regions by exact name (291) or a reviewed correspondence table (75); the region boundaries were unions of DHS boundary polygons for the same country where these exist, and admin-1 polygons otherwise (for example, Guinea-Bissau and the Central African Republic), and were used to locate MICS records in the Snow et al.\ areas as for DHS. As for DHS, all 13 area-level covariates were computed from the MICS microdata using the DHS definitions; education, recorded as level and grade, was converted to completed years, and facility delivery refers to the last birth in the two years before interview.


\emph{P.\ falciparum} prevalence estimates were taken from the Malaria Atlas Project (MAP) modelled surfaces of the parasite rate standardised to ages 2 to 10 years (\textit{Pf}PR$_{2-10}$), for 2000 - 2024 \cite{map,map2025}. Values for each Snow et al.\ area (below) were means of each annual surface over the area's polygon, weighted by the Gridded Population of the World version 4 (GPWv4) 2020 population density and grid-cell area \cite{gpwv4}; for a union of polygons, the population-weighted mean over its polygons was used.

\textit{Geographic units.} Records were located in the 520 malaria-endemic sub-national polygons used by Snow et al.\ to summarise \textit{P.\ falciparum} prevalence in sub-Saharan Africa \cite{snow:17}. Each record was located by the smallest administrative unit reported in its survey: the survey region, or, where a survey region spans several polygons, a finer location variable recorded in the survey file (for example state, province or district), matched to the polygons by boundary geometry, name or a reviewed alias. A record's unit was the polygon holding at least 90\% of the 2020 GPWv4 population of that location; otherwise it was the union of the polygons each holding at least 5\% of that population. Records in locations more than half of whose population lay in polygons classified as malaria free or as islands were excluded, as were records from the Comoros, São Tomé and Príncipe and Zanzibar, which have no endemic polygon. Of the 1,452 survey-by-area units in the primary analysis, 215 are unions of polygons (median 2, maximum 9 polygons), and the units of one survey can overlap; these unions hold 14\% of the mothers (Figure \ref{fig:map}). Rwanda has a single endemic polygon in the shapefile, so each Rwandan survey forms one unit. Figure \ref{fig:boundary_agreement} compares the survey regions with the Snow et al.\ polygons.

For each child trajectory, each of their recorded age bands (from birth up until interview or death) was assigned a \textit{Pf}PR$_{2-10}$ estimate based on the area of the mother's residence and the calendar year of entry into that age-band. For an age band to be eligible to be included in the primary analysis, the date of entry into the age band had to be within 60 months of the interview and it had to be completed by the time of the interview.
For national burden calculations, country-level prevalence was calculated using population-weighted means, with weights derived from the Gridded Population of the World version 4 population density surface for 2020 and grid-cell areas \cite{gpwv4}.


National age-specific all-cause death counts and rates were obtained from the IHME Global Burden of Disease data exports and used to calculate malaria-attributable mortality \cite{ihme}.
Comparator estimates of under-five malaria deaths were obtained from IHME and the CA-CODE 2026 series disseminated through the UN Inter-agency Group for Child Mortality Estimation (UN IGME) portal \cite{Perin:26,Villavicencio:24}. In addition we used IHME estimates for state specific all-cause mortality by age band in Nigeria (only country for which subnational estimates are available).

National child HIV incidence estimates for ages 0-14 years, expressed as new infections per 1,000 uninfected population, were obtained from the UNAIDS 2025 estimates distributed by UNICEF.
Corresponding adolescent incidence estimates for ages 15-19 years informed the imputation of missing child incidence. National GDP per capita, health expenditure per capita and political stability indicators were obtained from the World Bank World Development Indicators and Worldwide Governance Indicators \cite{wdi,wgi}.


\subsection*{Covariate adjustment}

Covariate selection was guided by the simplified causal processes shown in Figure~\ref{fig:dag}. This causal model was used to guide the selection of variables that are common causes of transmission intensity and mortality, that lie off the causal pathway. These are primarily variables that are proxies for socio-economic status and general health-system quality. We also selected variables that are determinants of childhood mortality from other causes. These include vaccination coverage, child HIV incidence and delivery in a health facility. We did not adjust for LLIN coverage (measured in DHS/MIS surveys) because this is a determinant of prevalence (LLINs reduce transmission) and corresponding reductions in mortality are mediated by reduction in transmission. For malnutrition, causality could go both ways and so measures of malnutrition could mediate the relationship between \textit{Pf}PR$_{2-10}$ and death, or confound it. To assess the potential impact of this, we carried out a sensitivity analysis with and without measures of malnutrition (wasting and stunting). 

Adjustment included 17 covariates: 13 area-level summaries (computed for each survey and area) and 4 national level measures. Area-level covariates were (means or percentages): maternal age at first birth (years); maternal education (years); household wealth-quintile score; urban versus rural residence; DTP3 coverage; measles vaccination coverage; proportion of live births delivered in a health facility; the percentage of non-first births with a preceding birth interval shorter than 24 months; household access to improved drinking water; access to improved sanitation; households with electricity; proportion of children with wasting; proportion with stunting. National covariates comprised child HIV incidence at ages 0–14 years; GDP per capita; health expenditure per capita; political stability, matched to the calendar year of entry into each age band. HIV incidence, GDP and health expenditure were natural-log transformed. All covariates were centred and scaled and entered linearly, with separate coefficients in each age-band model. We note that area-level summaries describe conditions at the time of the survey, not values in the past 60 months.

DHS summaries used the DHS women's sampling weights for Births Recode items and the household weights for Household Member and Height and Weight Recode items. MICS summaries were computed from the microdata with the MICS women's, household and children's sampling weights. Area-level summaries were calculated before complete-case selection; for DHS surveys they follow the DHS tabulation rules, with items not asked left out of the denominator and answers recorded as missing or don't know counted as `no'.

Missing area-level DTP3 and measles coverage was replaced by the corresponding national UNICEF estimate for the survey year (\url{https://immunizationdata.who.int/}). Remaining missing area-level covariates were replaced by the mean of the available values for the survey's other areas, weighted by the number of records in each area. South Sudan became independent from Sudan in July 2011; for band-entry years before 2011, its GDP per capita, health expenditure per capita and political stability were taken from Sudan's series for the same year, which admits the 2010 South Sudan MICS (57,509 child age-band records and 955 deaths, band-entry years 2005–2010). South Sudan's own child HIV incidence series was used throughout. Missing child HIV incidence was imputed using a model informed by adolescent incidence and temporal and geographic patterns; the primary analysis used fixed posterior-median imputations. Records with remaining missing required covariates were excluded. Imputation uncertainty was not propagated in the primary analysis. São Tomé and Príncipe has no child HIV incidence series; its surveys are excluded in any case, as it has no Snow et al.\ polygon.

%National child HIV incidence (new infections at ages 0--14 years per 1,000 uninfected population, both sexes) was obtained from the UNAIDS 2025 estimates distributed by UNICEF \cite{unaids2025}; the corresponding adolescent incidence at ages 15--19 years informed the imputation of missing child incidence. The UNICEF workbook has no incidence series for Liberia, so its child incidence was calculated from the UNAIDS 2026 estimates of annual new HIV infections at ages 0--14 years published on AIDSinfo \cite{aidsinfo}, as $1{,}000\times$ new infections divided by the under-five person-years implied by the IHME all-cause death count and rate, less the children aged 0--4 years living with HIV. Applied to the other countries, this formula reproduces the published child rates (median ratio of derived to published rate 0.99 across 761 country-years with at least 200 new infections, in 34 countries). National GDP per capita, health expenditure per capita and political stability indicators were obtained from World Bank data series.


\begin{figure}[htbp]\centering
\includegraphics[width=0.9\linewidth]{figures/Causal_diagram.pdf}
\caption{Assumed causal structure relating malaria transmission intensity (the malaria inoculation rate, proxied by \textit{Pf}PR$_{2-10}$) to all-cause child mortality, used to guide covariate selection. Solid arrows denote direct causes and dashed arrows effect modifiers. LLIN: long-lasting insecticide treated bednets; SES: socio-economic status.}
\label{fig:dag}
\end{figure}


\subsection*{Statistical model and model selection}

The unit of observation was the child age-band record. Following previous analyses of DHS/MIS birth histories \cite{Wakefield:19,Burstein:18}, each child's trajectory was divided into seven age bands as above. A band was included if the child entered it alive, entry fell within the 60 months before the interview, the entry year lay between 2000 and 2024, and the full band could have been completed by the interview date; this inclusion criteria applied equally to children who died, so a death in a band cut short by the interview was not counted (e.g. if a child had died a month before the interview at age 4.5 years they would not contribute to the 4-5 years age band, but only the previous age bands). The outcome was death versus survival during the age band. A child surviving to five years therefore contributed up to seven records; a child dying at eight months contributed three. Models were fitted to the child age-band records; because the exposure and all covariates are defined at the area and calendar-year level, this is equivalent to a binomial likelihood for the number of deaths among children entering each age band in each survey-by-area unit and entry month.

We fitted seven separate models, one per age band, because both the effect of \textit{Pf}PR$_{2-10}$ and the effects of the confounding variables could differ by age band. Each model was a binomial generalised additive model with a complementary log-log link and an offset for the band width in years, so that the linear predictor is the log annualised mortality hazard. The linear predictor comprised a penalised cubic regression spline of area \textit{Pf}PR$_{2-10}$ at time of band entry (basis dimension 5, with knots fixed at those of an earlier survey-region fit, \textit{Pf}PR$_{2-10}$ 0.02–94.8\%), a spline of calendar year at entry (basis dimension 6), the 17 standardised covariates entered linearly, and independent Gaussian random intercepts for survey, country and survey-by-area unit. Fewer than 200 records per age band (under 0.02\%) had a prevalence outside the knots, where the spline extrapolates linearly. Models were fitted by fast REML with discretised covariates using \texttt{bam} in the R package \texttt{mgcv} (version 1.9), with degrees of freedom penalisation parameter $\gamma = 2$ to force smoother curves. Hazard ratios between two prevalence levels $X$ and $Y$ were calculated as $\exp\{f_g(Y) - f_g(X)\}$ from the fitted spline $f_g$ of age band $g$. Confidence intervals used the coefficient covariance conditional on the estimated smoothing parameters and the single HIV imputation. Full model equations are given in the Supplement (SText~1).

\subsection*{Malaria-attributable mortality}

For each country and year, the malaria-attributable fraction of all-cause mortality for each age band $g$ was calculated as $1 - \exp\{f_g(0) - f_g(P)\}$, where $P$ is the population-weighted national mean \textit{Pf}PR$_{2-10}$ for that year, and applied to the IHME all-cause death counts and rates for the corresponding age group; the fitted curves were applied to all countries with MAP coverage, including those not contributing surveys to the model. IHME age groups were mapped to the model bands by combining early and late neonatal deaths (0-27 days) for the $<$1 month band, using the 1-5 and 6-11 month groups directly,  and assuming the 2-4 year group's deaths were spread equally across the three years. Estimates were produced annually for 2000–2024, in the 42 countries with national MAP estimates. National totals are point estimates: because the seven age-band models were fitted separately, their sampling covariance is not estimated and age-specific confidence intervals were not summed. For comparison with IHME and UN IGME (CA-CODE) malaria death estimates, deaths from each source were summed across the same 42 countries and expressed per 1,000 child-years using the under-five person-time implied by the IHME all-cause count and rate. The same calculation was repeated for the 36 Nigerian states and the Federal Capital Territory in 2024 using state-level IHME all-cause deaths by age and state MAP prevalence, with the state sum reconciled against the national estimate. For Figure \ref{fig:cumulative_risk}, the probability of dying before age 5 was calculated as a period life-table (synthetic cohort) probability, ${}_5q_0 = 1 - \exp(-\sum_g H_g)$, with each band's cumulative hazard $H_g$ from the IHME all-cause death rates; the probability caused by malaria is ${}_5q_0$ minus the probability with each band's hazard multiplied by the hazard ratio for \textit{Pf}PR$_{2-10}$ = 0 versus the country's current prevalence, and the IHME equivalent is ${}_5q_0$ minus the probability with IHME's malaria death rates removed.


\subsection*{Sensitivity analyses}
We carried out a suite of sensitivity analyses to assess whether restricting the data by seasonality of transmission, geographic sub-region or period, removing the nutrition covariates, or imputing missing covariates changed the estimated relationship between prevalence and age-specific all-cause mortality.
We refitted the 7 age-band models in subsets of the primary sample: Sahelian survey regions (boundary centroid at or north of 12°N and west of 36°E, excluding the Horn of Africa; 31 surveys in 8 countries with 18,768 deaths recorded); each UN M49 sub-region of Africa except Southern Africa, namely Western Africa (61 surveys in 15 countries, 54,881 deaths), Middle (Central) Africa (18 surveys in 7 countries, 17,314 deaths) and Eastern Africa (50 surveys in 12 countries, 31,858 deaths), while Southern Africa (7 surveys in Namibia, Eswatini and South Africa, 1,061 deaths) was not refitted because its \textit{Pf}PR$_{2-10}$ never exceeds 10.3\%; surveys conducted up to and including, versus after, the median survey year of 2014 (74 and 62 surveys, respectively); and child age bands entered up to and including, versus after, the median calendar year of entry (2012.8; 3,858,211 and 3,781,232 records). Subset fits used the same specification with knots at quantiles of the subset's own predictor values. In a second sensitivity analysis we imputed every remaining covariate gap rather than excluding records: whole-survey gaps in area-level indicators by chained-equation multiple imputation (predictive mean matching, 10 imputations), the missing 2001 political-stability round by interpolation, and health expenditure for Zimbabwe 2000–2009, Somalia 2000–2012 and all countries in 2024 from a generalised additive model on the observed panel, with draws. As in the primary analysis, South Sudan's GDP per capita, health expenditure per capita and political stability before 2011 were Sudan's values, held fixed. Child HIV incidence came from the extended incidence model (with a latent adolescent series for countries without one), with one posterior draw of each imputed or censored value per imputed dataset. This retained all 8,630,306 records with a Snow et al.\ area (161 surveys in 38 countries with 122,092 deaths); imputation uncertainty was propagated by refitting each of the ten imputed datasets with the smoothing parameters fixed at those of the point fit, and pooling with Rubin's rules.

% =====================================================================
\section*{Contributors}
XXXXX All authors approved the final version.

\section*{Declaration of interests}
We declare no competing interests.

\section*{Data sharing}
All input data are publicly available: DHS and MIS microdata from the DHS Program and MICS microdata from UNICEF (both after registration), MAP \textit{Pf}PR$_{2-10}$ surfaces, GPWv4 population, the Snow et al.\ sub-national polygons (published with the original article \cite{snow:17}), DHS survey-region boundaries from the DHS Spatial Data Repository, national boundaries from GADM, IHME estimates via the GBD client portal, UN IGME and CA-CODE estimates via the UN IGME portal, WHO/UNICEF immunisation coverage estimates, UNAIDS/UNICEF HIV estimates and World Bank indicators. Lake and reservoir outlines in Figure \ref{fig:map} are from Natural Earth (\url{https://www.naturalearthdata.com/}; public domain) and are used for display only. Analysis code and derived data tables are available at \url{https://github.com/Infectious-Diseases-Data-Observatory/malaria-mortality-U5}.


% =====================================================================
\bibliographystyle{vancouver}
\bibliography{ref}


% =====================================================================
\clearpage
\section*{Supplementary}
\renewcommand{\thefigure}{S\arabic{figure}}
\setcounter{figure}{0}

\renewcommand{\thetable}{S\arabic{table}}
\setcounter{table}{0}

\begin{figure}[htbp]\centering
\includegraphics[width=\linewidth]{Supplementary Figures/sfig_survey_map_and_timing.png}
\caption{DHS/MIS (circles) and MICS (triangles) surveys with complete birth histories conducted in sub-Saharan Africa since 2000: 124 DHS and MIS surveys and 47 MICS surveys in 41 countries. Filled symbols are the 136 surveys (99 DHS, 37 MICS) in 37 countries included in the primary analysis. Open symbols are the 35 excluded surveys: the five surveys from Lesotho (malaria free); the five from the Comoros (two) and São Tomé and Príncipe (three), which have no Snow et al.\ malaria-endemic polygon; and 25 surveys (14 DHS, all five MIS and 6 MICS) in which a required covariate was unavailable in every survey-by-area unit. Countries are grouped by UN sub-region (Central Africa: UN Middle Africa), and map shading shows the number of included surveys in each country. $n$ is the number of mothers (complete birth histories) contributing at least one child to the primary analysis (1,125,610 in total). Survey year is the year of fieldwork. The 46 MICS surveys without a complete birth history are not shown. DRC: Democratic Republic of the Congo; RC: Republic of the Congo; CAF: Central African Republic.}
\label{fig:survey_timing}
\end{figure}

\begin{figure}[htbp]\centering
\includegraphics[width=\linewidth]{Supplementary Figures/sfig_study_flow.png}
\caption{Sample inclusion for the primary analysis. Counts below the survey registry are child age-band records, not children. Records were placed in the Snow et al.\ malaria-endemic polygons through their smallest reported administrative unit (see Methods). External data are MAP \textit{Pf}PR$_{2-10}$ (population-weighted over each unit), WHO/UNICEF vaccine coverage estimates, UNAIDS child HIV incidence (new infections per 1,000 uninfected children aged 0–14 years) and World Bank GDP per capita, health expenditure per capita and political stability. For South Sudan, Sudan's values of the last three were used for band-entry years before 2011 (57,509 records).}
\label{fig:flow}
\end{figure}

\begin{figure}[htbp]\centering
\includegraphics[width=0.95\linewidth]{Supplementary Figures/sfig_pfpr_mortality_by_age.png}
\caption{Malaria prevalence and all-cause age-specific child mortality. Curves show the log hazard ratio of mortality in each age band relative to a prevalence of 20\%, from the seven age-band models of the primary analysis; lines are point estimates and shaded areas 95\% confidence intervals, drawn over the central 95\% of each band's exposure distribution.}
\label{fig:main_result}
\end{figure}


\begin{figure}
    \centering
    \includegraphics[width=0.95\linewidth]{Supplementary Figures/sfig_pfpr_splines_by_subgroup.png}
    \caption{Sensitivity analysis: Sahel. Log mortality hazard ratios relative to \textit{Pf}PR$_{2-10}$ = 20\% by age band, from the primary analysis (black, with 95\% confidence interval shaded) and refitted to records whose survey-region boundary centroid lies at or north of 12°N and west of 36°E, excluding the Horn of Africa (31 surveys in 8 countries, 18,768 deaths). The refit uses the primary specification with knots at quantiles of the subset's own values. Curves are drawn over the central 95\% of each sample's exposure. The subset is part of the primary sample, so the comparison is descriptive.}
    \label{fig:subgroups}
\end{figure}

\begin{figure}
    \centering
    \includegraphics[width=0.95\linewidth]{Supplementary Figures/sfig_pfpr_splines_by_subregion.png}
    \caption{Sensitivity analysis: UN M49 sub-regions. As Figure \ref{fig:subgroups}, refitted separately in Western Africa (61 surveys in 15 countries, 54,881 deaths), Middle (Central) Africa (18 surveys in 7 countries, 17,314 deaths) and Eastern Africa (50 surveys in 12 countries, 31,858 deaths). Southern Africa (7 surveys in 3 countries, 1,061 deaths) was not refitted because \textit{Pf}PR$_{2-10}$ never exceeds 10.3\%.}
    \label{fig:subgroups_region}
\end{figure}

\begin{figure}
    \centering
    \includegraphics[width=0.95\linewidth]{Supplementary Figures/sfig_pfpr_splines_by_period.png}
    \caption{Sensitivity analysis: period. As Figure \ref{fig:subgroups}, refitted separately in surveys conducted up to (74 surveys) or after (62 surveys) the median survey year, 2014, and in child age bands entered up to (3,858,211 records) or after (3,781,232 records) the median calendar year of entry, 2012.8.}
    \label{fig:subgroups_period}
\end{figure}

\begin{figure}
    \centering
    \includegraphics[width=0.95\linewidth]{Supplementary Figures/sfig_survival_vs_pfpr_by_subregion.png}
    \caption{Survival through each age band against \textit{Pf}PR$_{2-10}$ in the primary analysis sample, by UN sub-region. Each point is one survey $\times$ area unit $\times$ band-entry year (size: children entering the band). Lines are unadjusted spline fits for each sub-region (binomial generalised additive model with a complementary log-log link and a penalised cubic regression spline of \textit{Pf}PR$_{2-10}$, basis dimension 5, no covariates), drawn over the central 95\% of the sub-region's exposure, with 95\% intervals shaded. Southern Africa, whose prevalence spans 0–10\%, is shown as points only. All panels share the same y-axis; points below 93\% (0.6\% of children) are drawn at the lower edge as open triangles. The fits are descriptive: unlike the age-band models, they are not adjusted for confounders, calendar time or survey and area effects.}
    \label{fig:survival_by_region}
\end{figure}

\begin{figure}
    \centering
    \includegraphics[width=0.95\linewidth]{Supplementary Figures/sfig_pfpr_splines_imputed_covariates.png}
    \caption{Sensitivity analysis: imputed covariates. Log mortality hazard ratios relative to \textit{Pf}PR$_{2-10}$ = 20\% by age band, from the complete-case primary analysis (7,639,443 records, 136 surveys) and from all records with a Snow et al.\ area after imputing every remaining covariate gap (8,630,306 records, 161 surveys in 38 countries; see Methods). Shaded areas are 95\% confidence intervals; curves are drawn over the central 95\% of each sample's exposure. Pooling ten imputed datasets with Rubin's rules changed the 40\%-to-20\% hazard ratios by at most 0.001.}
    \label{fig:MI}
\end{figure}

\begin{figure}
    \centering
    \includegraphics[width=0.95\linewidth]{Supplementary Figures/sfig_pfpr_splines_without_nutrition.png}
    \caption{Sensitivity analysis: without the nutrition covariates. Log mortality hazard ratios relative to \textit{Pf}PR$_{2-10}$ = 20\% by age band, from the primary analysis (17 covariates) and from a refit on the same records without area-level wasting and stunting (15 covariates). Shaded areas are 95\% confidence intervals. The largest change in the hazard ratio for a reduction in prevalence from 40\% to 20\% is 0.010, and for 20\% to 0\% 0.036 (both at 1–5 months).}
    \label{fig:without_malnutrition}
\end{figure}


\begin{figure}[htbp]\centering
\includegraphics[width=0.95\linewidth]{Supplementary Figures/sfig_malaria_share_vs_pfpr.png}
\caption{Share of under-five deaths caused by malaria in 2024 against national population-weighted MAP \textit{Pf}PR$_{2-10}$ for 42 sub-Saharan African countries. The PfPR-ACM model share is attributable (deaths that would not occur at zero prevalence, including deaths caused indirectly by malaria); the IHME and UN IGME (CA-CODE) shares count deaths assigned to malaria as the cause.}
\label{fig:malaria_share_pfpr}
\end{figure}

\begin{figure}[htbp]\centering
\includegraphics[width=\linewidth]{Supplementary Figures/sfig_snow_boundary_agreement_map.png}
\caption{Survey regions of the DHS and MICS files against the Snow et al.\ polygons \cite{snow:17}, for the most recent included survey in each country and the Comoros MICS 2022, which is not analysed. Survey regions are classified by GPW 2020 population: same area (at least 90\% of the region's population lies in one Snow et al.\ polygon, which in turn has at least 90\% of its population in the region); region inside a larger polygon; polygons nested in the region (the region spans several polygons, each lying within it); and boundaries crossing. Regions lying mostly in malaria-free or island polygons, and regions with no Snow et al.\ polygon (the Comoros and Zanzibar), are shown separately; white areas are regions of a mapped survey that are not in the analysis sample. White lines: survey-region boundaries; thin dark lines: Snow et al.\ polygon boundaries; black lines: national borders. Countries without an included survey are light grey.}
\label{fig:boundary_agreement}
\end{figure}


\begin{table}[tbp]
\centering
\caption{Age-specific associations between malaria prevalence and all-cause mortality.}
\label{tab:primary-age-band-effects}
\small
\setlength{\tabcolsep}{3pt}
\renewcommand{\arraystretch}{1.15}
\begin{tabular*}{\linewidth}{@{\extracolsep{\fill}}lrrrcc@{}}
\toprule
\shortstack[l]{Age\\(months)} & \shortstack{Child-band\\records} & \shortstack{Deaths\\(\% of U5 deaths)} & \shortstack{PfPR\\EDF} & \shortstack{HR (95\% interval)\\40\% $\to$ 20\%} & \shortstack{HR (95\% interval)\\20\% $\to$ 0\%} \\
\midrule
$<1$ & 1,346,771 & 38,595 (36.7\%) & 2.42 & 0.94 (0.91--0.97) & 0.96 (0.91--1.01) \\
1--5 & 1,211,050 & 16,076 (15.3\%) & 2.32 & 0.92 (0.88--0.96) & 0.87 (0.81--0.93) \\
6--11 & 1,164,495 & 14,524 (13.8\%) & 2.61 & 0.84 (0.81--0.88) & 0.75 (0.69--0.82) \\
12--23 & 1,007,386 & 13,530 (12.9\%) & 3.61 & 0.88 (0.83--0.94) & 0.53 (0.47--0.61) \\
24--35 & 999,421 & 11,620 (11.1\%) & 3.73 & 0.89 (0.84--0.95) & 0.45 (0.40--0.52) \\
36--47 & 967,271 & 6,826 (6.5\%) & 3.37 & 0.89 (0.82--0.95) & 0.57 (0.49--0.66) \\
48--59 & 943,049 & 3,943 (3.7\%) & 3.32 & 0.94 (0.86--1.03) & 0.56 (0.48--0.66) \\
\midrule
Total & 7,639,443 & 105,114 (100.0\%) & -- & -- & -- \\
\bottomrule
\end{tabular*}
\par\vspace{0.5em}
\begin{minipage}{\linewidth}
\footnotesize
Deaths are observed deaths in the primary analysis sample; percentages use all 105,114 under-five deaths as the denominator.
Displayed percentages use largest-remainder rounding to one decimal place and sum to 100.0\%.
Records are child--age-band observations; a child can contribute multiple records.
EDF: effective degrees of freedom of the PfPR spline. HR: adjusted mortality hazard ratio for reducing PfPR from the first value to the second.
Models were fitted separately by age band using MAP prevalence and $\gamma=2$.
Intervals condition on fitted smoothing parameters, exposure, the fixed HIV imputation and filled area-level covariates; zero PfPR is below observed exposure support in every band.
\end{minipage}
\end{table}

\end{document}
